\documentclass[10pt,a4paper]{amsart}
\usepackage[margin=19mm]{geometry}
\usepackage{amsmath,amssymb,mathtools}
\usepackage[scr=rsfs,cal=euler]{mathalfa}
\usepackage{xfrac}
\usepackage[unicode,hidelinks,bookmarks=false]{hyperref}
\newcommand{\ie}{i.e.\ }
\newcommand{\cf}{cf.\ }
\newcommand{\resp}{respectively\ }
\newcommand{\Subsection}{\subsection}
\providecommand{\nobreakdash}{\nobreak\hbox{-}}
\setlength{\emergencystretch}{2em}
\multlinegap0pt
\usepackage{amssymb}
\usepackage{mathtools}
\usepackage[shortlabels]{enumitem}
\usepackage{xcolor}
\newtheorem{lemma}{Lemma}
\newcommand{\one}{\mathbf 1}
\title{Eigenvalues and eigenfunctions of the fractional Laplacian on the interval}
\author{Cheng Zhang}
\date{Source: arXiv:2608.23457v2 (CC BY 4.0)}
\begin{document}
\maketitle

\noindent\emph{Source and licence.} Eigenvalues and eigenfunctions of the fractional Laplacian on the interval. Version arXiv:2608.23457v2 (CC BY 4.0), distributed by arXiv under the Creative Commons Attribution 4.0 International licence, http://creativecommons.org/licenses/by/4.0/, as recorded in the arXiv licence metadata for this exact version and shown on the abstract page https://arxiv.org/abs/2608.23457v2. Full licence notice: \texttt{licenses/arxiv-2608.23457-CC-BY-4.0.txt}.\par\smallskip

\noindent\textit{Editorial scope.} Counted unit: the article's lemma lem:small-expansions (source0.tex lines 1877--1904). The article's complete original proof of it is delivered with it (source0.tex lines 1906--2247, one \texttt{proof} environment of 342 lines). No mathematics is written, repaired or re-derived: the statement and the proof are the article's own lines, in the article's own notation, and no retained line is substituted or re-flowed.  Retained uncounted as the source-local context the proof needs: lines 315--320; lines 323--325; lines 347--350.  Nothing was omitted from any retained span, so the package carries no omission record.  No retained line contains a citation, so the wrapper carries no bibliography and no bibliography entry.  No retained line contains a figure environment, an image-inclusion command or drawing code, so no image is delivered and the package declares no asset dependency.  Editorial formatting only, in addition: an AMS \texttt{amsart} preamble that restates the article's own declarations and macros used by the retained lines, namely the environments \texttt{lemma} on separate counters, together with the standard packages those lines need.  The retained environments print in this document as Lemma 1 in order of appearance, whereas the article prints the same items with the numbering of its own full document.  The complete native source of the article as distributed in the arXiv e-print for this exact version is bundled unchanged as \texttt{sources/208412/source0.tex}.
\begin{equation}\label{eq:F-zero}
 F_\alpha(t)\sim
 \frac{t^{\alpha/2}}
 {\sqrt{\alpha/2}\,\Gamma(\alpha/2)},
 \qquad t\downarrow0.
\end{equation}

\begin{equation}\label{eq:Gmass}
 \int_0^\infty G_\alpha(t)\,dt=M_\alpha.
\end{equation}

\begin{equation}\label{eq:Gsmall}
 G_\alpha(t)=K_\alpha t^{-1-\alpha}
 +K_{2,\alpha}t^{-1-2\alpha}+o_\alpha(t^{-1-2\alpha}),
\end{equation}

\begin{lemma}\label{lem:small-expansions}
Put
\[
 Q_\alpha=\int_0^\infty t^{-\alpha}F_\alpha(t)\,dt.
\]
The integral defining $Q_\alpha$ converges.  Moreover, as $n\to\infty$,
\begin{equation}\label{eq:FH-expansion}
\begin{aligned}
 &\int_0^{2\mu_n}F_\alpha(\rho)H_n(\rho)\,d\rho\\
 &\quad=\frac{K_\alpha Q_\alpha}{\alpha}(2\mu_n)^{-1-\alpha}
 +\frac{K_{2,\alpha}Q_\alpha}{\alpha}(2\mu_n)^{-1-2\alpha}\\
 &\qquad-\frac{K_\alpha\sqrt{\alpha/2}\,
 \Gamma(1-\alpha)\Gamma(1+2\alpha)}
 {\alpha\Gamma(1+\alpha)}(2\mu_n)^{-1-2\alpha}\\
 &\qquad+(-1)^{n+1}\frac{K_\alpha\cos\theta}{1+2\alpha}
 (2\mu_n)^{-1-2\alpha}
 +o_\alpha(\mu_n^{-1-2\alpha})
\end{aligned}
\end{equation}
and
\begin{equation}\label{eq:GH-expansion}
 \begin{aligned}
 \int_0^{2\mu_n}G_\alpha(2\mu_n-\rho)H_n(\rho)\,d\rho=\frac{K_\alpha M_\alpha}{1+2\alpha}
 (2\mu_n)^{-1-2\alpha}
 +o_\alpha(\mu_n^{-1-2\alpha}).
 \end{aligned}
\end{equation}
\end{lemma}

\begin{proof}
By \eqref{eq:F-zero},
\[
 t^{-\alpha}F_\alpha(t)=O_\alpha(t^{-\alpha/2}),
 \qquad 0<t\le1.
\]
Since $0<\alpha<1$, this is integrable at zero.  For $t\ge1$,
\eqref{eq:Gsmall} gives
\[
 t^{-\alpha}G_\alpha(t)=O_\alpha(t^{-1-2\alpha}).
\]
For $1\le A<B$, integration by parts gives
\[
 \left|\int_A^B t^{-\alpha}\sin(t+\theta)\,dt\right|
 \le3A^{-\alpha}.
\]
Thus $\int_1^\infty t^{-\alpha}\sin(t+\theta)\,dt$ converges, and
$Q_\alpha$ is well defined.

The expansion \eqref{eq:Gsmall}, its differentiated form, and complete
monotonicity imply
\begin{equation}\label{eq:G-tail-bounds-small}
 0\le G_\alpha(t)\le C_\alpha t^{-1-\alpha},
 \qquad
 0\le-G_\alpha'(t)\le C_\alpha t^{-2-\alpha},
 \qquad t\ge1.
\end{equation}
For $0<\rho<2\mu_n$,
\begin{equation}\label{eq:H-decomposition}
 \begin{aligned}
 H_n(\rho)
 =\frac{G_\alpha(2\mu_n)}{\alpha\rho^\alpha}+\int_0^\infty
 \frac{G_\alpha(2\mu_n+u)-G_\alpha(2\mu_n)}
 {(u+\rho)^{1+\alpha}}\,du.
 \end{aligned}
\end{equation}
Indeed,
$\int_0^\infty(u+\rho)^{-1-\alpha}\,du=\rho^{-\alpha}/\alpha$.
Moreover,
\[
 \begin{aligned}
 &\int_0^{2\mu_n}|F_\alpha(\rho)|\int_0^\infty
 \frac{|G_\alpha(2\mu_n+u)-G_\alpha(2\mu_n)|}
 {(u+\rho)^{1+\alpha}}\,du\,d\rho\\
 &\qquad\le\frac{G_\alpha(2\mu_n)}\alpha
 \int_0^{2\mu_n}\rho^{-\alpha}|F_\alpha(\rho)|\,d\rho<\infty.
 \end{aligned}
\]
Thus Fubini's theorem may be applied to \eqref{eq:H-decomposition}.  Using
\[
 (u+\rho)^{-1-\alpha}
 =\frac1{\Gamma(1+\alpha)}
 \int_0^\infty s^\alpha e^{-s(u+\rho)}\,ds,
\]
define
\[
 \Phi_n(s)=\int_0^{2\mu_n}e^{-s\rho}F_\alpha(\rho)\,d\rho,
 \qquad
 \Psi_n(s)=\int_0^\infty e^{-su}G_\alpha(2\mu_n+u)\,du.
\]
Then
\begin{equation}\label{eq:FH-Laplace-decomposition}
 \begin{aligned}
 &\int_0^{2\mu_n}F_\alpha(\rho)H_n(\rho)\,d\rho\\
 &\quad=\frac{G_\alpha(2\mu_n)}\alpha
 \int_0^{2\mu_n}\rho^{-\alpha}F_\alpha(\rho)\,d\rho\\
 &\qquad+\frac1{\Gamma(1+\alpha)}\int_0^\infty
 s^\alpha\Phi_n(s)
 \left(\Psi_n(s)-\frac{G_\alpha(2\mu_n)}s\right)ds.
 \end{aligned}
\end{equation}

We first expand the first term in \eqref{eq:FH-Laplace-decomposition}.
Since $2\mu_n+2\theta=n\pi$,
\[
 \cos(2\mu_n+\theta)=\cos(n\pi-\theta)
 =(-1)^n\cos\theta.
\]
Integration by parts gives
\[
 \begin{aligned}
 \int_{2\mu_n}^\infty
 \rho^{-\alpha}\sin(\rho+\theta)\,d\rho&=(2\mu_n)^{-\alpha}\cos(2\mu_n+\theta)
 -\alpha\int_{2\mu_n}^\infty\rho^{-1-\alpha}
 \cos(\rho+\theta)\,d\rho\\
 &=(-1)^n\cos\theta\,(2\mu_n)^{-\alpha}
 +O_\alpha(\mu_n^{-1-\alpha}).
 \end{aligned}
\]
By \eqref{eq:G-tail-bounds-small},
\[
 \begin{aligned}
 0&\le\int_{2\mu_n}^\infty
 \rho^{-\alpha}G_\alpha(\rho)\,d\rho\\
 &\le C_\alpha\int_{2\mu_n}^\infty
 \rho^{-1-2\alpha}\,d\rho
 =O_\alpha(\mu_n^{-2\alpha})
 =o_\alpha(\mu_n^{-\alpha}).
 \end{aligned}
\]
Consequently,
\begin{equation}\label{eq:weighted-F-tail}
 \int_0^{2\mu_n}\rho^{-\alpha}F_\alpha(\rho)\,d\rho
 =Q_\alpha+(-1)^{n+1}\cos\theta\,(2\mu_n)^{-\alpha}
 +o_\alpha(\mu_n^{-\alpha}).
\end{equation}
Equations \eqref{eq:Gsmall} and \eqref{eq:weighted-F-tail} yield
\begin{equation}\label{eq:FH-first-part}
\begin{aligned}
 &\frac{G_\alpha(2\mu_n)}\alpha
 \int_0^{2\mu_n}\rho^{-\alpha}F_\alpha(\rho)\,d\rho\\
 &\quad=\frac{K_\alpha Q_\alpha}{\alpha}(2\mu_n)^{-1-\alpha}+\left(
 \frac{K_{2,\alpha}Q_\alpha}{\alpha}
 +(-1)^{n+1}\frac{K_\alpha\cos\theta}{\alpha}
 \right)(2\mu_n)^{-1-2\alpha}
 +o_\alpha(\mu_n^{-1-2\alpha}).
\end{aligned}
\end{equation}

We next evaluate the second term in \eqref{eq:FH-Laplace-decomposition}.
Direct integration gives
\begin{equation}\label{eq:damped-sine}
 \begin{aligned}
 \int_0^{2\mu_n}e^{-s\rho}\sin(\rho+\theta)\,d\rho
 =\frac{s\sin\theta+\cos\theta
 -e^{-2\mu_ns}\bigl(s\sin(2\mu_n+\theta)
 +\cos(2\mu_n+\theta)\bigr)}{1+s^2}.
 \end{aligned}
\end{equation}
Since $G_\alpha\in L^1(0,\infty)$, \eqref{eq:damped-sine} and
\eqref{eq:Gmass} imply
\begin{equation}\label{eq:Phi-bounds-limit}
 \sup_{\substack{n\ge1\\s>0}}|\Phi_n(s)|\le C_\alpha,
 \qquad
 \Phi_n\left(\frac z{2\mu_n}\right)
 -\sqrt{\frac\alpha2}
 -(-1)^{n+1}\cos\theta\,e^{-z}\longrightarrow0
 \quad(z>0).
\end{equation}
Indeed, the sine integral in \eqref{eq:damped-sine}, with
$s=z/(2\mu_n)$, differs from
$\cos\theta+(-1)^{n+1}\cos\theta e^{-z}$ by $o_\alpha(1)$, while
\[
 \int_0^{2\mu_n}e^{-z\rho/(2\mu_n)}G_\alpha(\rho)\,d\rho
 \longrightarrow\int_0^\infty G_\alpha(\rho)\,d\rho=M_\alpha.
\]

Changing variables $u=2\mu_nv$ gives, for each $z>0$,
\begin{equation}\label{eq:Psi-limit}
 \begin{aligned}
 &(2\mu_n)^\alpha\left(
 \Psi_n\left(\frac z{2\mu_n}\right)
 -\frac{2\mu_nG_\alpha(2\mu_n)}z\right)\\
 &\quad=(2\mu_n)^{1+\alpha}\int_0^\infty e^{-zv}
 \bigl(G_\alpha(2\mu_n(1+v))-G_\alpha(2\mu_n)\bigr)\,dv\\
 &\quad\longrightarrow K_\alpha\int_0^\infty
 \bigl((1+v)^{-1-\alpha}-1\bigr)e^{-zv}\,dv.
 \end{aligned}
\end{equation}
Indeed, for $n\ge1$ and $v\ge0$, \eqref{eq:G-tail-bounds-small} gives
\[
 (2\mu_n)^{1+\alpha}
 \left|G_\alpha(2\mu_n(1+v))-G_\alpha(2\mu_n)\right|
 \le C_\alpha\bigl((1+v)^{-1-\alpha}+1\bigr).
\]
After multiplication by $e^{-zv}$, the right-hand side is integrable in
$v$.  Thus \eqref{eq:Gsmall} and dominated convergence prove
\eqref{eq:Psi-limit}.  We next prove two bounds for its left-hand side.
First,
\[
 \begin{aligned}
 &(2\mu_n)^\alpha\left|
 \Psi_n\left(\frac z{2\mu_n}\right)
 -\frac{2\mu_nG_\alpha(2\mu_n)}z\right|\\
 &\quad\le(2\mu_n)^\alpha\left(
 \Psi_n\left(\frac z{2\mu_n}\right)
 +\frac{2\mu_nG_\alpha(2\mu_n)}z\right)\\
 &\quad\le\frac{2(2\mu_n)^{1+\alpha}G_\alpha(2\mu_n)}z
 \le\frac{C_\alpha}{z}.
 \end{aligned}
\]
Second, integration by parts in $v$ gives
\[
 \Psi_n\left(\frac z{2\mu_n}\right)
 -\frac{2\mu_nG_\alpha(2\mu_n)}z
 =\frac{(2\mu_n)^2}{z}\int_0^\infty
 e^{-zv}G_\alpha'(2\mu_n(1+v))\,dv,
\]
and hence
\[
 \begin{aligned}
 &(2\mu_n)^\alpha\left|
 \Psi_n\left(\frac z{2\mu_n}\right)
 -\frac{2\mu_nG_\alpha(2\mu_n)}z\right|\\
 &\qquad\le\frac{C_\alpha}{z}
 \int_0^\infty e^{-zv}(1+v)^{-2-\alpha}\,dv
 \le\frac{C_\alpha}{z^2}.
 \end{aligned}
\]
Therefore
\begin{equation}\label{eq:Psi-dominating-bound}
 (2\mu_n)^\alpha\left|
 \Psi_n\left(\frac z{2\mu_n}\right)
 -\frac{2\mu_nG_\alpha(2\mu_n)}z\right|
 \le C_\alpha\min\{z^{-1},z^{-2}\}.
\end{equation}
Because $0<\alpha<1$,
$z^\alpha\min\{z^{-1},z^{-2}\}$ is integrable on $(0,\infty)$.
After the change of variables $s=z/(2\mu_n)$, equations
\eqref{eq:Phi-bounds-limit}--\eqref{eq:Psi-dominating-bound} give
\begin{equation}\label{eq:FH-scaled-limit}
 \begin{aligned}
 &(2\mu_n)^{1+2\alpha}
 \frac1{\Gamma(1+\alpha)}\int_0^\infty
 s^\alpha\Phi_n(s)
 \left(\Psi_n(s)-\frac{G_\alpha(2\mu_n)}s\right)ds\\
 &\quad=\frac{K_\alpha}{\Gamma(1+\alpha)}
 \int_0^\infty z^\alpha
 \left(\sqrt{\frac\alpha2}
 +(-1)^{n+1}\cos\theta\,e^{-z}\right)\\
 &\qquad\qquad\times\int_0^\infty
 \bigl((1+v)^{-1-\alpha}-1\bigr)e^{-zv}\,dv\,dz
 +o_\alpha(1).
 \end{aligned}
\end{equation}
The two integrals obtained by reversing the order of integration are
\begin{align}
 \int_0^\infty\bigl((1+v)^{-1-\alpha}-1\bigr)v^{-1-\alpha}\,dv
 &=-\frac{\Gamma(1-\alpha)\Gamma(1+2\alpha)}
 {\alpha\Gamma(1+\alpha)},
 \label{eq:first-constant-integral}\\
 \int_0^\infty\bigl((1+v)^{-1-\alpha}-1\bigr)
 (1+v)^{-1-\alpha}\,dv
 &=\frac1{1+2\alpha}-\frac1\alpha.
 \label{eq:second-constant-integral}
\end{align}
Using
\[
 \int_0^\infty z^\alpha e^{-zv}\,dz
 =\Gamma(1+\alpha)v^{-1-\alpha},
 \qquad
 \int_0^\infty z^\alpha e^{-z(1+v)}\,dz
 =\Gamma(1+\alpha)(1+v)^{-1-\alpha},
\]
equations \eqref{eq:FH-scaled-limit}--\eqref{eq:second-constant-integral}
show that the second term in \eqref{eq:FH-Laplace-decomposition} equals
\begin{equation}\label{eq:FH-second-part}
 \begin{aligned}
 &\left[-\frac{K_\alpha\sqrt{\alpha/2}\,
 \Gamma(1-\alpha)\Gamma(1+2\alpha)}
 {\alpha\Gamma(1+\alpha)}
 +(-1)^{n+1}K_\alpha\cos\theta
 \left(\frac1{1+2\alpha}-\frac1\alpha\right)\right]\\
 &\qquad\times(2\mu_n)^{-1-2\alpha}
 +o_\alpha(\mu_n^{-1-2\alpha}).
 \end{aligned}
\end{equation}
Adding \eqref{eq:FH-first-part} and \eqref{eq:FH-second-part} proves
\eqref{eq:FH-expansion}.

For \eqref{eq:GH-expansion}, the substitutions
$x=2\mu_n-\rho$ and $y=2\mu_n+u$ give
\begin{equation}\label{eq:GH-double-integral}
 \begin{aligned}
 &\int_0^{2\mu_n}G_\alpha(2\mu_n-\rho)H_n(\rho)\,d\rho\\
 &\qquad=\int_0^{2\mu_n}G_\alpha(x)
 \int_{2\mu_n}^\infty
 \frac{G_\alpha(y)}{(y-x)^{1+\alpha}}\,dy\,dx.
 \end{aligned}
\end{equation}
For every fixed $x\ge0$, the change of variables $y=2\mu_nv$ gives
\[
 \begin{aligned}
 (2\mu_n)^{1+2\alpha}\int_{2\mu_n}^\infty
 \frac{G_\alpha(y)}{(y-x)^{1+\alpha}}\,dy=(2\mu_n)^{1+\alpha}\int_1^\infty
 G_\alpha(2\mu_nv)
 \left(v-\frac{x}{2\mu_n}\right)^{-1-\alpha}\,dv.
 \end{aligned}
\]
For all sufficiently large $n$ and $v\ge1$,
\[
 (2\mu_n)^{1+\alpha}G_\alpha(2\mu_nv)
 \left(v-\frac{x}{2\mu_n}\right)^{-1-\alpha}
 \le C_\alpha v^{-2-2\alpha}.
\]
Since the left-hand side converges to
$K_\alpha v^{-2-2\alpha}$, dominated convergence gives
\begin{equation}\label{eq:GH-inner-limit}
 (2\mu_n)^{1+2\alpha}\int_{2\mu_n}^\infty
 \frac{G_\alpha(y)}{(y-x)^{1+\alpha}}\,dy
 \longrightarrow\frac{K_\alpha}{1+2\alpha}.
\end{equation}
If $0\le x\le\mu_n$, then $y-x\ge y/2$ for $y\ge2\mu_n$, and
\begin{equation}\label{eq:GH-inner-bound}
 \begin{aligned}
 (2\mu_n)^{1+2\alpha}\int_{2\mu_n}^\infty
 \frac{G_\alpha(y)}{(y-x)^{1+\alpha}}\,dy\le C_\alpha(2\mu_n)^{1+2\alpha}
 \int_{2\mu_n}^\infty y^{-2-2\alpha}\,dy
 \le C_\alpha.
 \end{aligned}
\end{equation}
For every $x\ge0$, equations \eqref{eq:GH-inner-limit} and
\eqref{eq:GH-inner-bound} give
\[
 \begin{aligned}
 0\le\one_{\{x\le\mu_n\}}G_\alpha(x)(2\mu_n)^{1+2\alpha}
 \int_{2\mu_n}^\infty
 \frac{G_\alpha(y)}{(y-x)^{1+\alpha}}\,dy\le C_\alpha G_\alpha(x),
 \end{aligned}
\]
and the left-hand side converges pointwise to
$K_\alpha G_\alpha(x)/(1+2\alpha)$.  Since
$G_\alpha\in L^1(0,\infty)$, dominated convergence gives
\begin{equation}\label{eq:GH-main-range}
 \begin{aligned}
 \int_0^{\mu_n}G_\alpha(x)\int_{2\mu_n}^\infty
 \frac{G_\alpha(y)}{(y-x)^{1+\alpha}}\,dy\,dx=\frac{K_\alpha M_\alpha}{1+2\alpha}
 (2\mu_n)^{-1-2\alpha}
 +o_\alpha(\mu_n^{-1-2\alpha}).
 \end{aligned}
\end{equation}
For $\mu_n<x<2\mu_n$, put $\rho=2\mu_n-x$ and $u=y-2\mu_n$.
By \eqref{eq:G-tail-bounds-small},
\[
 \begin{aligned}
 &\int_{\mu_n}^{2\mu_n}G_\alpha(x)\int_{2\mu_n}^\infty
 \frac{G_\alpha(y)}{(y-x)^{1+\alpha}}\,dy\,dx\\
 &\quad=\int_0^{\mu_n}G_\alpha(2\mu_n-\rho)
 \int_0^\infty\frac{G_\alpha(2\mu_n+u)}
 {(\rho+u)^{1+\alpha}}\,du\,d\rho\\
 &\quad\le C_\alpha(2\mu_n)^{-2-2\alpha}
 \int_0^{\mu_n}\int_0^\infty
 (\rho+u)^{-1-\alpha}\,du\,d\rho\\
 &\quad=\frac{C_\alpha}{\alpha}(2\mu_n)^{-2-2\alpha}
 \int_0^{\mu_n}\rho^{-\alpha}\,d\rho\\
 &\quad=O_\alpha(\mu_n^{-1-3\alpha})
 =o_\alpha(\mu_n^{-1-2\alpha}).
 \end{aligned}
\]
Combining this estimate with \eqref{eq:GH-double-integral} and
\eqref{eq:GH-main-range} proves \eqref{eq:GH-expansion}.
\end{proof}

\end{document}
